% !TeX root = adelic-line-bundle-0916.tex
\section{adelic-line-bundle-0916}

\subsection[status=draft,tag=0000000,label=\labelNFM{subsec:height-and-volumes}]{Height and volumes}

Recall the projective case.
Let $K$ be a number field, $\calX/\calO_K$ projective variety, and $\bar{\calL} \in \aPic(\calX)$.
Then $H^0(\calX,\calL)$ is a $\calO_K$-module with a hermitan metric $\|l\|_{v,sup} = \sup_{x \in \calX_v(\bbC)} \|l(x)\|_v$.
Define
\[ \hat{H}^0(X,\bar{\calL}) = \{l \in H^0(X,\calL), \|l\| \leq 1  \}. \]
it is a finite set.
Define
\[ \hat{h}^0(X,\bar{\calL}) = \log \# \hat{H}^0(X,\bar{\calL}) \]
We say that $\bar{\calL}$ is effective if $\hat{h}^0(X,\bar{\calL}) > 0$.
If $\bar{\calL}$ has an effective section $s$, $\overline{D} := \adiv(s)$, then $x \in X(\bar{K}) \setminus \Supp D, s|_X \neq 0, \adeg \calL|_{\bar{x}} \geq 0$ and hence $h_{\bar{\calL}}(x) \geq 0$.

Define $\calX(\bar{\calL}) = - \log \vol (H^0(\calX,\calL)_\bbR / H^0(\calX,\calL))$.
The $\vol$ is the Haar measure on $\hat{H}^0(\calX,\calL)_\bbR$ has measure 1.

\begin{theorem}[label=\labelNFM{thm:hilbert-samuel}, status=draft, tag=0000000]
	Let $\calL_K$ be ample, $\calL \in \aPic(\calX/\calO_K)_\sm$, $\bar{\calL}$ relative semiample.
	Then, for all $\bar{\calM} \in \aPic(\calX/\calO_K)$,
	\[ \chi(n\bar{\calL} + \bar{\calM}) = \frac{n^d}{d!} \bar{\calL}^d + o(n^d). \]
\end{theorem}

We have \emph{arithmetic ampleness}.
Let $\calX/\calO_K$ be a projective arithmetic variety and $\bar{\calL} \in \aPic(\calX)_\sm$ with $\calL_K$ being ample.
We say that $\bar{\calL}$ is arithmetically ample if for any $\bar{\calM} \in \aPic(\calX)$, the $H^0(\calX,n\calL+\calM)$ is generated by $\hat{H}^0(\calX,n\bar{\calL}+\bar{\calM})$ for $n \gg 0$.

We say that $\bar{\calL}$ is arithmetically positive if $\cal_K$ is ample, $\bar{\calL}$ is relatvie semipositive, and there exist $c > 0$ such that $h_{\bar{\calL}}(x) \geq c$ for all $x \in X(\bar{K})$.

\begin{theorem}[label=\labelNFM{thm:cre}, status=draft, tag=0000000]
	Assume $\calL_K$ is ample and $\bar{\calL}$ is relative semipositive. TFAE
	\begin{enumerate}
		\item $\bar{\calL}$ is arithmetically ample.
		\item $\bar{\calL}$ is arithmetically positive.
		\item $\forall Y \subset \calX_K$, $h_{\bar{\calL}}(Y) \geq 0$.
	\end{enumerate}
\end{theorem}

Under these condition above,
\[ \hat{h}^0(n \bar{\calL} + \bar{\calM}) = \frac{n^d}{d!} \bar{\calL}^d + o(n^d) \]

\subsection[status=draft,tag=0000000,label=\labelNFM{subsec:essmin}]{Essential minima}

Let $\calX,\bar{\calL} \in \aPic(\calX)$, define
\[ e_i(\bar{\calL}) \coloneqq \sup_{\codim Y = i} \inf_{x \in (X\setminus Y)(\bar{K})} h_{\bar{\calL}}(x), \quad i = 1,\cdots,d. \]
If $\bar{\calL}$ is effective, then $e_1(\bar{\calL}) \geq 0$.

\begin{theorem}[label=\labelNFM{thm:fundomental-inequality}, status=draft, tag=0000000]
	Assume $\calL_K$ ample, $\bar{\calL}$ relative semipositive.
	Then
	\[ e_1(\bar{\calL}) \geq h_{\bar{\calL}}(\calX) \geq \frac{e_1(\bar{\calL})+\cdots+e_d(\bar{\calL})}{d}. \]
	When $e_1(\bar{\calL}) = \cdots = e_d(\bar{\calL})$, we have the Equidistribution theorem.
\end{theorem}


\subsection[status=draft,tag=0000000,label=\labelNFM{subsec:big}]{Arithmetic bigness}

\begin{theorem}[label=\labelNFM{thm:yuan}, status=draft, tag=0000000]
	Let $\calX/\calO_K$ of dimension $d$, $\bar{\calL} = \bar{\calL}_1 - \bar{\calL}_2$ arithmetically positive.
	Then for all $\bar{\calM}$, $\chi(n\bar{\calL}+\bar{\calM}) \geq n^d/d! (\bar{\calL}_1^d - d \bar{\calL}_1^{d-1} \bar{\calL}_2) - o(d)$.
\end{theorem}

Define
\[ \hat{\vol}(\bar{\calL}) = \limsup_{n \to \infty} \frac{\hat{h}^0(\calX,n\bar{\calL})}{n^d/d!}. \]
the suplim indeed a limit.

If $\hat{\vol}(\bar{\calL}) > 0$, we say that $\bar{\calL}$ is big.
We say that $\bar{\calL}$ is pseudo-effective if for any big $\overline{M}$, $\bar{\calL} + \bar{\calM}$ is big.

\begin{corollary}[label=\labelNFM{cor:I}, status=draft, tag=0000000]
	If $\bar{\calL}$ is nef, then $\hat{\vol}(\bar{\calL}) = \bar{\calL}^d$.

\end{corollary}

The volume is continuous.


\subsection[status=draft,tag=0000000,label=\labelNFM{subsec:quasi}]{Quasi projective case}

An adelic divisor $\bar{D} \in \aDiv(\calU/\bbZ)$ is called effective if $\bar{D}$ can be represented by a seqeuence of effective model divisors.













































